% Transcription — fonds Alexandre Grothendieck, shelfmark 5, batch 3 (pages 41-52).
% Established from the facsimile archives/batches/5/batch-03.pdf (a mirror of
% https://grothendieck.umontpellier.fr/5.pdf), per the skill
% transcribe-grothendieck. The batch file carries no cover sheet: PDF page N is
% archivists' page N + 40. Checked against the pencilled numbers bottom-left
% ("41" on the first leaf, "42", "43", "44", "45", "46", "47", "48", "49", "50",
% "51", "52" on the following ones), which confirm the alignment end to end.
%
% Pass: Fable 5.1 (claude-fable-5-1), 2026-09-19 — first pass, unchecked against the pages by a human.
%
% What the batch is. Twelve leaves, of three kinds, and two leaves that are
% nobody's mathematics. Page 42 is one leaf in his hand, in blue ink, closing a
% note on homogeneous polynomial maps and their polarisations that began in
% batch 2 (the run is not readable from here; a \note records the boundary).
% Pages 44 and 46 are two leaves in his hand, in black ink, headed « Cristaux »
% by him: the conditions a system of elements π^{(n)} ∈ πA must satisfy for
% x ↦ x^{(n)} = π^{(n)} y^n (x = πy) to be a divided-power structure on the
% principal ideal J = πA, with three examples — A torsion-free over Z, A a
% Z_(p)-algebra with π = p, and A a discrete valuation ring of absolute
% ramification index e, where the condition v(π) ≥ e·sup_n v_p(n!)/n = e/(p−1)
% gives, for π a uniformiser, the boxed inequality e ≤ p − 1. Pages 45 and 47
% are the versos of 44 and 46: two pages, numbered 23 and 22, of a typescript
% of EGA IV § 18.13 (the "completed" module of differentials Ω̄¹_{B/A} of a
% topological algebra, the representability (18.13.4.1) of continuous
% derivations, and Proposition (18.13.5)), gone over in his hand — insertions,
% a cancelled clause, two marginal notes, a marginal formula. Following the
% convention of folders 4, 48 and 55, the typed text is given in full where
% his hand touches it, and here it touches nearly every line, so the two
% pages are given entire; the typist's own overstrikes (x's over a mistyped
% word) are not recorded, what he struck by hand is. The typescript reads
% 22 (leaf 47) then 23 (leaf 45): page 22 ends "co-" and page 23 begins
% "ïncide". The pages are kept in the archivists' order, and notes say so.
% Page 48 is an otherwise blank leaf carrying, in his hand, « lettre à
% Deligne / sur coh. p-adique »: the cover of the letter that follows, and it
% gets no \page; its words are the section heading. Pages 49 to 52 are that
% letter, a typed carbon dated by hand « 10.12.1965 » — the « lettre (1965) »
% the inventory announces — to Deligne: a simplification of the proof of the
% duality theorem avoiding relative purity, a relative form of the theorem
% over a base with a sheaf of rings A killed by an integer n prime to the
% residue characteristics, and then the conjecture that the transcendental
% comparison between de Rham and ℓ-adic cohomology over C should have a
% p-adic analogue over a complete non-archimedean field, through Tate's
% rigid-analytic space X^an, a rigid-analytic Poincaré lemma, and, in
% characteristic p, the map H*(X, F_p) → H*_DR(X) whose composite with the
% Frobenius is the only thing it can be. The carbon has no arrows; they are
% restored from the text and a note says so once.
%
% Skipped, without description: page 41, a typescript (a course announcement
% and a seminar programme) scanned upside down, nothing in his hand; page 43,
% a typescript on Banach categories and Spin(p,q), nothing in his hand.
%
% Notation fixed for the batch. Divided powers are written as he writes them,
% π^{(n)}, x^{(n)}; his circled 1 labelling the three conditions is set (1).
% Z is \mathbf{Z}. In the EGA typescript the underlined r (the radical) is
% \mathfrak{r} as in the printed EGA, the overlined Ω is \overline{\Omega},
% and « Dér » is \mathrm{D\acute{e}r}. In the letter the typescript's
% underlined sheaf symbols (O_X, H^q, Ω) are kept underlined, as typed. The
% one tikz-cd diagram, (18.13.5.1), has three vertical arrows and two rows
% of three, transcribed as typed.
\documentclass[11pt,a4paper]{article}
% Le préambule commun à toutes les transcriptions du fonds Grothendieck.
%
% Il tient en une idée : **l'incertitude se compose, elle ne se résout pas.**
% Une transcription qui lisse un mot douteux a détruit la seule chose pour
% laquelle on la faisait — un lecteur doit pouvoir savoir, à chaque endroit,
% ce qui était sur le feuillet et ce qui a été ajouté. Les six macros qui
% suivent sont donc l'appareil critique, et non de la décoration.
%
% Les mêmes commandes sont reconnues par `scripts/render.mjs`, qui produit la
% vue de lecture du site. Une macro ajoutée ici doit l'être aussi là-bas,
% faute de quoi le rendu échouera bruyamment — ce qui est voulu : un rendu
% silencieusement faux est pire qu'un rendu absent.

\ProvidesPackage{grothendieck}[2026/08/08 Transcriptions du fonds Grothendieck]

\usepackage{amsmath,amssymb,amsthm}
\usepackage{mathrsfs}
\usepackage{stmaryrd}
% Les diagrammes commutatifs sont partout dans ce fonds : c'est la langue
% même de la géométrie algébrique de Grothendieck, et une transcription qui
% les rendrait en prose ne transcrirait rien.
\usepackage{tikz-cd}
% Français par défaut — c'est la langue des feuillets — mais l'anglais est
% chargé aussi : la traduction et le résumé sont des documents anglais, et la
% typographie française (espaces avant les deux-points) y serait une faute.
% Ces documents basculent par \selectlanguage{english} en tête de corps.
\usepackage[english,french]{babel}
\usepackage{fontspec}
\usepackage{geometry}
\geometry{a4paper,margin=25mm,marginparwidth=28mm}
\usepackage{marginnote}
\usepackage{xcolor}
% ulem, not soul. `soul`'s \st and \ul are built on a character-by-character
% scan that XeLaTeX's font machinery breaks: the document fails with an
% undefined control sequence rather than degrading. ulem does the same two jobs
% and is Unicode-engine safe. `normalem` stops it hijacking \emph, which the
% transcriptions use for Grothendieck's own emphasis.
\usepackage[normalem]{ulem}
\usepackage{hyperref}
\hypersetup{colorlinks=true,linkcolor=black,urlcolor=black,pdfborder={0 0 0}}

% --- Métadonnées du lot -----------------------------------------------------
% Reprises telles quelles par le script de rendu, qui les affiche en tête de
% la vue de lecture. Elles fixent ce que le fichier prétend couvrir : sans
% elles, une transcription flotte sans point d'attache dans un fonds de seize
% mille feuillets.

\newcommand{\folder}[1]{\gdef\grothFolder{#1}}
\newcommand{\batch}[1]{\gdef\grothBatch{#1}}
\newcommand{\leaves}[2]{\gdef\grothFirst{#1}\gdef\grothLast{#2}}
\newcommand{\foldertitle}[1]{\gdef\grothFoldertitle{#1}}
\gdef\grothFolder{?}\gdef\grothBatch{?}\gdef\grothFirst{?}\gdef\grothLast{?}\gdef\grothFoldertitle{}

% --- Le résumé --------------------------------------------------------------
%
% Il ouvre la lecture modernisée, sous le titre « Résumé », et il est composé
% en petit. Ce n'est pas un ornement : c'est le seul passage du document qui
% ne soit pas des mathématiques, et un lecteur doit pouvoir le reconnaître
% avant de l'avoir lu — puis le sauter s'il connaît déjà le sujet, ou s'y
% arrêter s'il ne le connaît pas. Le filet qui le suit marque la reprise du
% corps.

\newenvironment{resume}
  {\small\setlength{\parskip}{0.45em}}
  {\par\medskip\hrule\medskip}

% --- Mots-clés ---------------------------------------------------------------
%
% En anglais, en clôture du résumé de la lecture modernisée : le vocabulaire
% moderne sous lequel chercher ce que les feuillets construisent. Le manifeste
% du site (`npm run manifest`) les extrait pour étiqueter la cote entière —
% cette ligne est la source unique des tags, il n'y a pas de fichier à côté.
\newcommand{\keywords}[1]{\par\smallskip\noindent
  {\footnotesize\textsc{Keywords} --- \textit{#1}}\par}

% --- L'appareil critique ----------------------------------------------------

% Le numéro de feuillet, en marge. C'est l'ancre qui permet de régler un
% désaccord sur une formule en regardant *un* feuillet plutôt qu'une chemise
% de six cents. Le site s'en sert aussi pour tourner les pages du fac-similé
% quand on fait défiler la transcription.
\newcommand{\leaf}[1]{%
  \marginnote{\footnotesize\color{gray!70}\textsf{#1}}%
  \label{leaf:#1}\ignorespaces}

% Illisible : on ne devine pas. Un mot inventé qui se lit comme les autres est
% le pire résultat possible.
\newcommand{\ill}{\textcolor{red!60!black}{[\,\ldots\,]}}

% Lecture douteuse : le mot est donné, et signalé comme incertain.
\newcommand{\uncertain}[1]{\uline{#1}}

% Ajout de l'éditeur — une lettre manquante, un mot sous-entendu.
\newcommand{\add}[1]{\textcolor{blue!50!black}{[#1]}}

% Biffé par Grothendieck lui-même. Ce qu'il a rayé fait partie du document :
% une rature dit souvent où la pensée a changé de direction.
\newcommand{\struck}[1]{\sout{#1}}

% Note du transcripteur, sur l'état matériel du feuillet.
\newcommand{\note}[1]{\par\smallskip{\small\color{gray!60!black}\textit{[#1]}}\par\smallskip}

% Note marginale de Grothendieck, distincte du corps du texte.
\newcommand{\marginal}[1]{\par\smallskip\noindent\hspace{1em}%
  {\small\color{gray!60!black}#1}\par\smallskip}

% --- Titre ------------------------------------------------------------------

\AtBeginDocument{%
  \begin{center}
    {\large\bfseries Fonds Alexandre Grothendieck --- cote n\textsuperscript{o}~\grothFolder}\\[2pt]
    {\itshape\grothFoldertitle}\\[4pt]
    {\small feuillets \grothFirst--\grothLast\ (lot \grothBatch)}\\[2pt]
    {\footnotesize\color{gray!60!black}Transcription établie d'après le fac-similé numérisé
     par l'Université de Montpellier.\\
     Les lectures incertaines sont soulignées, les illisibles notées [\,\ldots\,],
     les ajouts entre crochets.}
  \end{center}
  \vspace{1em}}

\endinput


\folder{5}
\batch{3}
\pages{41}{52}
\dating{1965-1969}
\watermark{Édition de démonstration}
\foldertitle{Cristaux et cohomologie de De Rham : généralités, correspondances Deligne et Berthelot (1965-1969) : notes manuscrites (s.d.), lettre (1965).}

\begin{document}

\section*{Applications polynomiales homogènes}

\note{titre de l'éditeur. Le feuillet 42 est la fin d'une note sur les
applications polynomiales commencée dans le lot précédent (page 40 et avant),
que ce lot ne peut pas lire ; les notations $\Phi$, $M$, $N$ y sont fixées.
Sur les premières lignes du feuillet, $\Phi$, $M$ et $N$ sont soulignés, puis
ne le sont plus ; « de degré $n$ » de la première ligne est en interligne.}

\page{42}

Une application polynomiale homogène de degré $n$ $\Phi : M \to N$ définit
par la formule
\[
\Phi\Bigl(\sum_i t_i x_i\Bigr)
= \sum_{p_1 + \cdots + p_r = n} \Phi_{p_1 \ldots p_r}(x_1, \ldots, x_r)\,
t_1^{p_1} \cdots t_r^{p_r}
\]
[les $t_i$ des indéterminées] des applications polynomiales
\[
\Phi_{p_1 \ldots p_r} : M \times \cdots \times M \to N
\]
multihomogènes de degrés $p_1, \ldots, p_r$, d'où des applications
\[
\varphi_{p_1, \ldots, p_r} : \underbrace{M \times \cdots \times M}_{r} \to N
\]
(définies si $r \geq 1$, $p_1, \ldots, p_r \geq 0$, $\sum p_i = n$). Ces
applications sont reliées par des relations
\[
\varphi_{p_1, \ldots, p_r}\Bigl(\sum_{1 \leq i \leq s_1} x_{1,i},\
\sum_{1 \leq i \leq s_2} x_{2,i},\ \ldots,\ \sum_{i \leq s_r} x_{r,i}\Bigr)
= \sum_{\substack{\sum_{1 \leq i_1 \leq s_1} q_{1,i_1} = p_1 \\ \sum_{1 \leq i_2 \leq s_2} q_{1,i_2} = p_2 \\ \cdots \\ \sum_{1 \leq i_r \leq s_r} q_{r,s_r} = p_r}}
\varphi_{q_{1,i_1}, \ldots, q_{r,s_r}}(x_{1,1}, \ldots, x_{r,s_r})
\]
\note{les conditions de sommation sont écrites sous une accolade, l'une sous
l'autre ; la seconde se lit $q_{1,i_2}$ sur la page, où l'on attend
$q_{2,i_2}$, et l'indice de la dernière est écrit $q_{r,s_r}$ tel quel}
\marginal{il suffit \ill{} pour \ill{} des \uncertain{arguments} \ill{}}
\note{note marginale de six courtes lignes, accolée à la relation, dont seuls
les premiers mots se lisent}

\uncertain{Plus} relations de symétrie \note{souligné}

On suppose que l'on a \ill{} \uncertain{vérifié} que … \note{la phrase
s'interrompt sur des points de suspension ; le reste du feuillet est blanc}

\section*{Cristaux}

\note{titre de sa main, souligné, en tête du feuillet 44. La note se poursuit
au feuillet 46 ; les feuillets 45 et 47, qui en sont les versos, portent deux
pages d'un tapuscrit annoté de sa main, données ci-dessous à leur place dans
la numérotation}

\page{44}

Soit $J$ un idéal principal d'un anneau $A$, $J = \pi A$. Formellement pour
y définir des puissances divisées, on est amené à définir des
\[
\frac{(\pi y)^n}{n!} = \frac{\pi^n}{n!}\, y^n = \pi^{(n)} y^n
\]
donc à définir des éléments
\[
\pi^{(n)} \in \pi A \qquad (n \geq 1)
\]
qui satisfassent aux conditions suivantes.

\struck{$\pi^{(m)} \pi^{(n)} = \frac{(m+n)!}{m!\,n!}\, \pi^{(m+n)}$ ; et si
$\pi z = 0$, \ill{} $\pi^{(n)} (y+z)^n = \pi^{(n)} \uncertain{z}^n$, ce qui
\ill{} \ill{} des $\pi^{(n)} y$} \note{trois lignes cernées d'un trait et
biffées, premier état des conditions (1) ; dans la seconde le dernier facteur
se lit $z^n$, où le sens attend $y^n$}

\[
(1)\qquad
\begin{cases}
\pi^{(1)} = \pi \\ \pi^{(m)} \pi^{(n)} = \dfrac{(m+n)!}{m!\,n!}\, \pi^{(m+n)} \\ \pi^{(p)} \Bigl(\dfrac{\pi^{(q)}}{\pi}\Bigr)^{p} = \dfrac{(pq)!}{p!\,(q!)^p}\, \pi^{(pq)}
\end{cases}
\]
\note{le « 1 » de l'étiquette est cerclé sur la page, ici et à chaque rappel}

NB. Comme $\pi^{(q)} \in \pi A$, $\pi^{(q)}/\pi$ est défini \uncertain{mod.
addition} d'un $z$ tel que $\pi z = 0$, mais comme $\pi^{(p)} \in \pi A$, est
déterminé sans ambiguïté l'élément $\pi^{(p)} \bigl(\pi^{(q)}/\pi\bigr)^p$.

Alors on trouve une structure de puissances divisées (sur $J = \pi A$)
\note{« sur $J = \pi A$ » en interligne, le $J$ tracé ici comme un I gras}
\[
x^{(n)} = (\pi y)^{(n)} = \pi^{(n)} y^n \qquad \text{pour } x = \pi y \in J,
\]
en notant que $x^{(n)}$ ne dépend pas du choix de $y$ t.q. $x = \pi y$.

\emph{Exemple 1.} Supposons que $A$ soit sans torsion sur $\mathbf{Z}$. Alors
les \struck{\ill{}} relations
\[
\pi^n = n!\, \pi^{(n)}
\]
déterminent uniquement les $\pi^{(n)}$ en fonction de $\pi$, et leur existence
équivaut au fait que
\[
\pi^n \in n!\, A \qquad \forall n \geq 2 .
\]
Les relations (1) sont satisfaites automatiquement.

\emph{Exemple 2.} Supposons $A_0 = \mathbf{Z}_{(p\mathbf{Z})}$, \struck{Alors}
$\pi = p$. Alors \note{« $\pi = p$. Alors », précédé d'un « Alors » biffé, est
écrit en bout de la ligne qui précède le titre « Exemple 2 » ; la virgule
après $\mathbf{Z}_{(p\mathbf{Z})}$ l'y rattache} les conditions précédentes
sont satisfaites, les $p^n/n!$ sont des entiers $p$-adiques. Donc si $A$ est
\ill{} dans laquelle tout nb premier $q \neq p$ est \uncertain{inversible}
(i.e. les \uncertain{caractéristiques} résiduelles sont $p$, $0$), \ill{} une
algèbre sur $A_0$, on trouve par spécialisation des $\pi^{(n)} \in A$,
$\pi^{(n)} = p^n/n!$, qui

\section*{Tapuscrit EGA IV, § 18.13, page 23 (annoté)}

\note{titre de l'éditeur. Copie d'un tapuscrit paginé « IV-977-23 », les
chiffres 977 et 23 à la main, au verso du feuillet 44. La page 22 du même
tapuscrit est au feuillet 47, verso du feuillet 46, et se lit avant celle-ci :
son dernier mot, « co- », se poursuit ici par « ïncide ». Les interventions à
la main sont données comme biffures, ajouts et notes marginales ; les ratures
propres de la dactylographie ne sont pas relevées. L'énoncé de la proposition
est souligné dans le tapuscrit}

\page{45}

ïncide avec le module de différentielles analytiques défini dans (18.13.1) ;
en effet \note{« en effet » en interligne, au-dessus d'un mot raturé à la
machine} (les définitions étant celles de (18.13.1)), il résulte du fait que
$\overline{\Omega}^1_{A/k}$ est de type fini et de (0, 20.4.8) que la
dérivation $\bar d_{A/k}$ se factorise en $A \xrightarrow{d_{A/k}}
\Omega^1_{A/k} \xrightarrow{u_0} \overline{\Omega}^1_{A/k}$ et que tout
$A$-homomorphisme $\Omega^1_{A/k} \xrightarrow{u} L$, où $L$ est un $A$-module
de type fini, se factorise nécessairement en $\Omega^1_{A/k}
\xrightarrow{u_0} \overline{\Omega}^1_{A/k} \xrightarrow{v} L$, si bien que
$\mathrm{Ker}(u)$ contient toujours $\mathrm{Ker}(u_0)$, d'où notre assertion.

On notera que si $\overline{\Omega}^1_{B/A}$ est un $B$-module de type fini,
il représente le foncteur covariant $L \rightsquigarrow
\mathrm{D\acute{e}r}_A(B, L)$ dans la catégorie des $B$-modules de type fini.
Il est clair que si $\Omega^1_{B/A}$ lui-même est de type fini (ce qui sera
par exemple le cas (0, 20.4.7) lorsque $B$ est une $A$-algèbre de type fini),
on a $\overline{\Omega}^1_{B/A} = \Omega^1_{B/A}$.

\textbf{Proposition (18.13.5).} --- Soient $A$ un anneau topologique,
\struck{\ill{}} B et C deux \note{« B et C deux » à la main, remplaçant un mot
biffé} $A$-algèbres topologiques qui \struck{\ill{}} sont des \note{« sont
des » en interligne} anneaux \struck{\ill{}} semi-locaux noethériens et dont
les topologies sont les topologies préadiques, $\rho : B \to C$ un
$A$-homomorphisme continu. Alors :

(i) Si \struck{\ill{}} $\overline{\Omega}^1_{C/A}$ est un $C$-module de type
fini \struck{\ill{}}, $\overline{\Omega}^1_{C/B}$ est un $C$-module de type
fini.

(ii) Si $\overline{\Omega}^1_{B/A}$ est un $B$-module de type fini et si $C$
est une $B$-algèbre finie, $\overline{\Omega}^1_{C/A}$ est un $C$-module de
type fini.

(iii) Supposons vérifiées, soit les hypothèses de (ii), soit \struck{\ill{}}
l'hypothèse de (i) et la condition supplémentaire que $\Omega^1_{B/A}$ est un
$B$-module de type fini ; alors il existe deux $C$-homomorphismes
\struck{déterminés de façon uniques} \marginal{style \ill{}} \note{la note
marginale, deux mots dont le second n'est pas lu, est accolée par un trait à
la biffure} : $\bar u : \overline{\Omega}^1_{C/A} \to
\overline{\Omega}^1_{C/B}$ et $\bar v : \overline{\Omega}^1_{B/A} \otimes_B C
\to \overline{\Omega}^1_{C/A}$ tels que le diagramme

\begin{tikzcd}[column sep=small, row sep=small, nodes={font=\scriptsize}]
\Omega^1_{B/A} \otimes_B C \arrow[r, "v"] \arrow[d, "r^{\prime} \otimes 1"'] & \Omega^1_{C/A} \arrow[r, "u"] \arrow[d, "r"'] & \Omega^1_{C/B} \arrow[r] \arrow[d, "r^{\prime\prime}"'] & 0 \\
\overline{\Omega}^1_{B/A} \otimes_B C \arrow[r, "\bar v"'] & \overline{\Omega}^1_{C/A} \arrow[r, "\bar u"'] & \overline{\Omega}^1_{C/B} \arrow[r] & 0
\end{tikzcd}

\note{numéroté (18.13.5.1) à gauche du diagramme ; la page s'arrête sur lui}

\section*{Cristaux (suite)}

\page{46}

satisfait aux conditions précédentes (1), donc une structure de puissances
divisées canonique sur $pA$.

\emph{Exemple 3.} Soit $A$ intègre, \struck{\ill{}} tel que les seules
\struck{\ill{}} \note{une ligne biffée, non lue} caractéristiques résiduelles
soient $p$ et $0$. Soit $\pi$ un élément tel que \struck{les}
\uncertain{caractéristiques} résiduelles de $A/\pi A$ soient $p$. À quelle
condition peut-on trouver des $\pi^{(n)}$ ? Supposons que $A$ soit un anneau
de valuation discrète. Il faut \struck{\ill{}} $\pi^n/n! \in A$, i.e.
\[
v(\pi^n/n!) \geq 0 \qquad \text{pour tout } n
\]
i.e.
\[
n\, v(\pi) - v(n!) \geq 0 \qquad \text{pour tout } n .
\]
Soit \struck{\ill{}} $v(p) = e$ l'indice de ramification absolu de $A$, donc
\struck{$v(n!)$} pour un entier $N$, $v(N) = e\, v_p(N)$. \struck{Soit donc
\ill{} $\rho = v(\pi)$.} La condition s'écrit \struck{\ill{}}
\[
v(\pi) \geq e\, \frac{v_p(n!)}{n} \qquad \text{pour tout } n
\]
i.e.
\[
v(\pi) \geq e\, \lambda(p) \qquad \text{où } \lambda(p) = \sup_n
\frac{v_p(n!)}{n} = \frac{1}{p-1}
\]
\marginal{\uncertain{Legendre}} \note{en marge droite, au-dessus de la formule
donnant $\lambda(p)$, qui est elle-même écrite en marge et rattachée au
texte par une accolade} est une constante ne dépendant que de $p$, et qui est
$\leq 1$ \marginal{\struck{\ill{} $\lambda(p)$}} … Donc on trouve la condition
\[
v(\pi) \geq \frac{e}{p-1} .
\]
Supposons que $\pi$ soit l'uniformisante, \uncertain{ou} $v(\pi) = 1$, on
trouve la condition
\[
\boxed{e \leq p-1}
\]

\section*{Tapuscrit EGA IV, § 18.13, page 22 (annoté)}

\note{titre de l'éditeur. Page « IV-977-22 » du tapuscrit dont la page 23 est
au feuillet 45 ci-dessus ; elle se lit avant celle-ci. Le $\mathfrak{r}$
souligné du tapuscrit est rendu $\mathfrak{r}$}

\page{47}

(18.13.4) Le résultat de (18.13.2), \struck{joint au fait que pour une
$k$-algèbre analytique $A$, $\overline{\Omega}^1_{A/k}$ soit un $A$-module de
type fini,} peut encore s'exprimer en disant que dans la catégorie des
$A$-modules de type fini, $\overline{\Omega}^1_{A/k}$ est un objet
représentant le foncteur covariant $L \rightsquigarrow
\mathrm{D\acute{e}r}_k(A, L)$ \note{« $(0_{\mathrm{III}}, 8.1.11)$ » ajouté à
la main au-dessus de la ligne, accolé à « représentant le foncteur
covariant »}. Nous allons généraliser ce résultat. Considérons d'abord un
anneau topologique $A$, et une $A$-algèbre topologique $B$ qui est un anneau
semi-local noethérien \marginal{à quoi sert \uncertain{l'hypothèse} \ill{} ?
Prendre \ill{} topologie plus \uncertain{générale} ?} \note{« semi-local
noethérien » cerné d'un trait, un astérisque renvoyant à la note marginale,
écrite en biais}, muni de sa topologie $\mathfrak{r}$-préadique
($\mathfrak{r}$ étant le radical de $B$). Alors, pour tout $B$-module de type
fini $L$, toute $A$-dérivation $D$ de $B$ dans $L$ est continue pour les
topologies $\mathfrak{r}$-préadiques sur $B$ et $L$, puisque l'on a
$D(\mathfrak{r}^{j+1}) \subset \mathfrak{r}^j L$ pour tout entier $j$ ; en
vertu de (0, 20.4.8) elle s'écrit donc $D = u \circ d_{B/A}$, où $u :
\Omega^1_{B/A} \to L$ est un homomorphisme de $B$-modules (nécessairement
continu pour les topologies $\mathfrak{r}$-préadiques, et on sait que dans le
cas considéré la topologie canonique \note{« canonique » en interligne} de
$\Omega^1_{B/A}$ est la topologie \struck{\ill{}} $\mathfrak{r}$-préadique
\note{« $\mathfrak{r}$-préadique » en interligne, au-dessus du mot biffé}
(0, 20.4.5)). Nous désignerons par $\overline{\Omega}^1_{B/A}$ le quotient de
$\Omega^1_{B/A}$ par l'intersection des noyaux des homomorphismes
$\Omega^1_{B/A} \to L$ pour les $B$-modules $L$ \note{« $L$ » ajouté en
interligne} de type fini (ou, ce qui revient au même puisque $L$ est
noethérien, l'intersection des sous-modules $M$ de $\Omega^1_{B/A}$ tels que
$\Omega^1_{B/A}/M$ soit de type fini). \marginal{On a donc
$\mathrm{Hom}_B(\overline{\Omega}^1_{B/A}, L) = \mathrm{Hom}_B(\Omega^1_{B/A},
L)$ pour tout $B$-module de type fini $L$.} \note{note marginale, insérée ici
par un trait} Nous désignerons par $\bar d_{B/A}$ ou simplement $\bar d$
l'application composée $B \xrightarrow{d_{B/A}} \Omega^1_{B/A}
\xrightarrow{r} \overline{\Omega}^1_{B/A}$ ; c'est une $A$-dérivation
\note{« $A$- » en interligne}, et $\overline{\Omega}^1_{B/A}$ est engendré par
les éléments $\bar d_{B/A}(x)$ où $x$ parcourt $B$, en vertu de (0, 20.4.7) ;
la définition précédente montre que l'application $v \rightsquigarrow v \circ
\bar d_{B/A}$ est un isomorphisme de $B$-modules, fonctoriel en $L$
\[
(18.13.4.1)\qquad \mathrm{Hom}_B(\overline{\Omega}^1_{B/A}, L)
\xrightarrow{\ \sim\ } \mathrm{D\acute{e}r}_A(B, L) =
\mathrm{D\acute{e}r}.\mathrm{cont}_A(B, L)
\]
\note{« $= \mathrm{D\acute{e}r}.\mathrm{cont}_A(B, L)$ » ajouté à la main}
$L$ étant toujours supposé de type fini. Lorsque l'on est dans les conditions
de (18.13.2), le $A$-module $\overline{\Omega}^1_{A/k}$ qui vient d'être
défini co- \note{la phrase se poursuit à la page 23 du tapuscrit, feuillet 45
ci-dessus}

\section*{Lettre à Deligne sur coh. p-adique}

\note{titre de sa main sur le feuillet 48, par ailleurs blanc, qui n'a pas de
numéro de page ici ; il coiffe la lettre dactylographiée des feuillets 49 à
52, la « lettre (1965) » de l'inventaire. Le double dactylographié n'a pas
les flèches, restituées d'après le texte, ni les crochets de décalage ; le
premier paragraphe, de « Je vous propose » à « constant $\mathbf{Z}/n\mathbf{Z}$ », est
délimité par deux crochets tracés à la main dans la marge}

\page{49}

10.12.1965 \note{date à la main, en haut à droite}

Cher Deligne,

Je vous propose une simplification pour la démonstration du théorème de
dualité, qui permet d'éviter tout recours au théorème de pureté relative.
Vous vous rappelez qu'on était réduit au cas où $f : X \to Y$ est de
dimension relative $1$, $F = A_X$ et $G = A_Y$. Le procédé \note{« proéédé »
dans la frappe} de passage à la limite de Exp VI, par. 6, permet de supposer
la base noethérienne, et même si on y tien\add{t} de dimension finie (car de
type fini sur $\mathbf{Z}$). On raisonne par récurrence sur $n = \dim X$. Si
$n = 0$, on sait le vérifier. Supposons le théorème démontré en dimension
$< n$. Se localisant sur $Y$, on peut supposer $Y$ strictement local. Alors,
si $y$ est son point fermé, on a $\dim(Y - y) \add{<} n$, et comme on est
réduit à prouver le théorème séparément pour $X \times_Y (Y - y)$ et $X
\times_Y y$, on gagne. --- Autre remarque : le théorème de dualité peut se
démontrer, essentiellement de la même façon et avec le même énoncé, pour un
morphisme lisse $f : X \to Y$ compactifiable, lorsque $Y$ est muni d'un
faisceau d'anneaux $A$ quelconque tel que il existe un entier $n$ premier aux
caractéristiques résiduelles annulant $A$, et $X$ d'un faisceau d'anneaux
$B$, et $f$ étant donné comme morphisme de $(X, B)$ dans $(Y, A)$, de sorte
qu'on a un homomorphisme de faisceau d'anneaux $f^{-1}(A) \to B$. On définit
alors
\[
f^{!}(K^{\bullet}) = R\,\underline{\mathrm{Hom}}^{\bullet}_{f^{-1}(A)}
\bigl(B,\ f^{-1}(K) \otimes T_{X/Y}\,\add{[}2d\add{]}\bigr) .
\]
La définition de l'homomorphisme trace et la démonstration du théorème de
dualité se décomposent alors en le cas où $f^{-1}(A) \to B$ est un
isomorphisme, qui se traite comme le cas $A = (\mathbf{Z}/n\mathbf{Z})_Y$, et
le cas où $f = \mathrm{id}$, qui est trivial. Je pense qu'il vaut le coup
d'inclure cette forme générale du théorème de dualité, soit de prime abord,
soit à la fin dans un numéro-\uncertain{rattrappage}. Je n'ai pas regardé si
par hasard il pourrait se déduire du cas particulier

\page{50}

$A = (\mathbf{Z}/n\mathbf{Z})_Y$ comme simple corollaire, mais ça
m'étonnerait. La même remarque s'applique d'ailleurs également au théorème de
dualité globale, qui s'énonce pour des faisceaux d'anneaux plus généraux que
le faisceau constant $\mathbf{Z}/n\mathbf{Z}$.

Il me semble que la relation de nature transcendante qui lie la cohomologie
de De-Rham ou de Hodge aux cohomologies $\ell$-adiques, lorsque le corps de
base est $\mathbf{C}$, via la cohomologie entière, devrait avoir un analogue
lorsque le corps de base est alg.\ clos \note{« alg clos » ajouté à la main}
valué complet non archimédien, à corps résiduel de caractéristique $p$, et
corps des fractions $K$ \note{« $K$ » ajouté à la main} de caractéristique
zéro, savoir que $H^{*}(X, \mathbf{Z}_p)$ doit s'envoyer alors de façon
$\mathbf{Z}_p$-linéaire dans la cohomologie de De Rham $H^{*}(X)$, qui serait
isomorphe alors à $H^{*}(X, \mathbf{Z}_p) \otimes_{\mathbf{Z}_p} K$. Pour
définir un tel homomorphisme, on pourrait penser à utiliser l'espace
rigide-analytique au sens de Tate défini par $X$ (NB c'est un site qui ne
peut être défini par une topologie au sens ordinaire, ce qui explique les
difficultés conceptuelles initiales de Tate pour arriver à définir la notion
d'espace rigide-analytique), soit $X^{\mathrm{an}}$. On espère que les
arguments habituels prouveront que pour des faisceaux de torsion, la
cohomologie de $X_{\text{ét}}$ coïncide avec celle de $X^{\mathrm{an}}$, et
un excès d'optimisme nous ferait espérer que la cohomologie de
$X^{\mathrm{an}}$ à coefficients constants $\mathbf{Z}_p$ est bien la limite
projective des cohomologies à coefficients dans les $\mathbf{Z}/p^n\mathbf{Z}$,
donc coïncide avec la cohomologie $p$-adique $H(X_{\text{ét}}, \mathbf{Z}_p)$
(définie précisément comme une limite analogue). Si cela était vrai,
l'immersion canonique de $\mathbf{Z}_p$ dans les fonctions constantes sur
$X^{\mathrm{an}}$ semblerait permettre d'envoyer $H(X^{\mathrm{an}},
\mathbf{Z}_p)$ dans la cohomologie de De Rham de $X^{\mathrm{an}}$, qui par
Gaga rigide-analytique n'est autre que celle de $X$. L'« excès » plus haut
serait justifié essentiellement par un « lemme

\page{51}

de Poincaré » rigide-analytique, disant que sur $X^{\mathrm{an}}$, le
complexe de De-Rham est bien une résolution du faisceau des fonctions
constantes à valeurs dans le corps de base. Notez bien que c'est certainement
trivial pour l'analytique $p$-adique ordinaire, étant une simple question de
séries converge\struck{a}ntes dans ce cas, mais que la question pour la
topologie rigide-analytique (qui est beaucoup plus grossière) est plus
délicate. Je vous avoue d'ailleurs que j'ai de grands doutes que les choses
marchent aussi simplement que ça, i.e. que la cohomologie rigide-analytique
à coefficients constants tout bêtes donne les « bons » nombres de Betti. Je
crois me rappeler que Tate m'avait dit que c'était déjà faux pour le $H^1$
des courbes elliptiques --- il faudrait lui demander quel était son
argument ; apparemment, on trouverait un peu plus que pour la cohomologie
étale, mais quand même pas toute la cohomologie. Il reste plausible
néanmoins, quoi qu'il en soit, que la topologie rigide-analytique aura son
rôle à jouer dans ces questions.

Notez qu'en caractéristique $p > 0$, on a un homomorphisme évident de
$H^{*}(X, \mathbf{F}_p)$ dans la cohomologie de De-Rham, comme on voit en
calculant cette dernière pour la cohomologie étale et non pour la cohomologie
de Zariski (ce qui donne le même résultat, puisque les composantes du
complexe de De Rham sont quasi-cohérents), et utilisant la suite spectrale en
cohomologie étale
\[
H^{*}(X) \Leftarrow H^p\bigl(X, \underline{H}^q(\underline{\Omega})\bigr) .
\]
Cet homomorphisme se factorise d'ailleurs à travers $H^{*}(X, \mathbf{F}_p)
\to H^{*}(X, \underline{O}_X)$, ce dernier s'envoyant dans $H^{*}(X)$ grace à
l'homomorphisme de puissance $p$-ème $f \rightsquigarrow f^p$, induisant un
isomorphisme $\underline{O}_X \simeq \underline{H}^0(\underline{\Omega})$. Le
composé des homomorphismes canoniques $H^n(X, \underline{O}_X) \to H^n(X) \to
H^n(X, \underline{O}_X)$ (ce dernier

\page{52}

résultant de l'autre suite spectrale pour la cohomologie de De Rham) ne peut
guère être autre chose que l'homomorphisme de Frobenius. Il faut dire que
tout ça est bien éculé, et qu'on ne pourra dire des choses vraiment
intéressantes et nouvelles qu'en faisant appel à la « vraie » cohomologie
$p$-adique.

Bien cordialement

\end{document}
